problem:  
Let \(N_{k,l} = \mathrm{SU}(3)/\mathrm{U}(1)_{k,l}\) be an Aloff–Wallach space with the standard \(\mathrm{Ad}(\mathrm{U}(1)_{k,l})\)-invariant decomposition  
\[
\mathfrak{m} = \mathfrak{m}_1 \oplus \mathfrak{m}_2 \oplus \mathfrak{m}_3,\qquad 
g_{\lambda_1,\lambda_2,\lambda_3} = \lambda_1^2 Q|_{\mathfrak{m}_1} + \lambda_2^2 Q|_{\mathfrak{m}_2} + \lambda_3^2 Q|_{\mathfrak{m}_3},
\]
where \(Q\) is the negative Killing form on \(\mathfrak{su}(3)\).  
Consider the cohomogeneity‑one manifold
\[
M_{k,l} = (0,\infty) \times N_{k,l}
\]
equipped with a \(\mathrm{SU}(3)\)-invariant metric of the form
\[
g = dr^2 + g_{\lambda_1(r),\lambda_2(r),\lambda_3(r)}.
\]
Suppose there exists an isotropy inclusion \( \mathrm{U}(1)_{k,l} \subset H \subset \mathrm{SU}(3)\) so that the metric can extend smoothly at \(r=0\) to a singular orbit of type \(H\).  
Possible examples include \(H = T^2\) (yielding the flag manifold \(\mathrm{SU}(3)/T^2\)) and, for special pairs \((k,l)\) where \(\mathrm{U}(1)_{k,l} \subset \mathrm{SU}(2)\), \(H = \mathrm{SU}(2)\) (yielding the 5-sphere \(\mathrm{SU}(3)/\mathrm{SU}(2) \cong S^5\)).

**Problem:**  
Determine whether there exist complete, noncompact, cohomogeneity‑one **gradient steady or expanding Ricci solitons**
\[
\mathrm{Ric}(g) + \nabla^2 f = 0 \quad (\text{steady}), \qquad
\mathrm{Ric}(g) + \nabla^2 f = \lambda g,\ \lambda < 0 \quad (\text{expanding}),
\]
on such \(M_{k,l}\), for some smooth potential \(f(r)\), satisfying the appropriate regularity conditions at any admissible singular orbit \(H\) and completeness as \(r\to\infty\).

Specifically:
1. Do there exist smooth ODE solutions \((\lambda_1(r),\lambda_2(r),\lambda_3(r),f(r))\) near \(r=0\) that meet the analytic initial conditions ensuring smooth extension over the corresponding singular orbit?  
2. For which isotropy types \(H\) (and associated pairs \((k,l)\)) can such complete soliton trajectories exist?  
3. Do any of these solitons exhibit **asymptotically conical behavior** with cone links given by \(N_{k,l}\) equipped with a homogeneous Einstein metric \(g_{\lambda_1^\infty,\lambda_2^\infty,\lambda_3^\infty}\)?  
4. How do asymptotic curvature decay rates and potential growth depend analytically on the isotropy parameters \((k,l)\)?  

Why is it a "good" problem:  
This formulation corrects the geometric assumptions on singular orbits and captures a genuinely open and profound question at the interface of **homogeneous geometry, Ricci soliton analysis, and dynamical systems**.  
The problem probes whether noncompact gradient steady or expanding Ricci solitons can exist with **Aloff–Wallach principal orbits**, extending known soliton constructions (such as the Bryant and Dancer–Wang solitons) to a new setting with richer isotropy structure.  

The system reduces to a delicate nonlinear ODE problem whose completeness and asymptotic conditions intertwine **Lie‑algebraic representation data** \((k,l)\) with **geometric PDE behavior**. Solving or even partially understanding this would yield insights into how **cohomogeneity‑one soliton dynamics** interact with **homogeneous Einstein geometries**, potentially revealing new families of solitons or geometric obstructions.  
The question is concise yet deeply open, with natural links among curvature flows, symmetry, and boundary regularity—fully meeting the “narrow entrance, profound depth” requirements of a good research-level mathematical problem.